\section{Q1：最大安全载荷与精确组批}\label{sec:q1}
\subsection{为什么额定载荷不足以判断安全}
额定载荷只限制装载质量，不保证山区往返能源充足。载货去程既要完成水平飞行，又要把整机与货物抬升至 DEM 决定的巡航高度；卸货后的返程仍有飞行和爬升开销。还须保留题设返航电量，因此应先回答每个服务区—机型组合“最多可安全运多少”，再讨论怎样把实际货箱装成批次。

\subsubsection{服务区—机型安全载荷边界}
Q1 限定单点往返 $O01\to i\to O01$。定义
\begin{equation}
q^*_{gi}=\max\{q\in[0,Q_g]:
E_{g,0i}(q)+E_{g,i0}(0)\le(1-\rho_g)E_g\}.\label{eq:qmax}
\end{equation}
该边界同时扣除了空载返程和储备，不能仅用额定载荷或单程航程判断。若空载往返已不可行，则该组合记为不可达；若额定载荷可行，则 $q^*_{gi}=Q_g$；其余情况求能量边界。不可达与“可达但只能承载零质量”是不同状态，输出中分别标记。

\subsubsection{单调性与二分依据}
当官方参数满足 $L_g^0>L_g^F>0$ 时，式\eqref{eq:range}对 $q$ 非增，故水平项 $E_gd/L_g(q)$ 非减，爬升项也随 $q$ 非减；空载返程为常量。因此载荷可行集是一个区间，二分能定位其上界。提高 $\rho_g$ 只会缩小能量预算，不能提高安全载荷，这是数值求解与边界测试共用的单调性依据。
\paperfig{fig03_payload_boundary}{额定载荷与能量共同形成服务区—机型能力边界。上图列出全部组合，下图给出安全载荷占额定值比例；C 型部分远端服务区先触及能源限制。各点对应不同地形，连线不能解释为控制其他因素后的纯距离响应。}{payload}
结果如图\ref{fig:payload}：C 型有 \CRestricted{} 个服务区不能满额定载荷往返，其最小安全载荷为 \CMinPayload{} kg；B 型最小为 \BMinPayload{} kg。A 型在基准储备下达到额定载荷。这表明机型“大”不等于任意服务区都能满载，能力判断必须同时考虑其能量与航程参数。

\subsection{不可拆货箱的集合划分模型}
安全载荷是连续质量边界，而需求由带有编号、质量和体积的不可拆货箱构成。即使总质量可运，也未必能按平均质量切成可行批次。对服务区 $i$ 的货箱集合 $\mathcal B_i$，枚举非空子集 $S$ 和机型 $g$，仅保留同时满足式\eqref{eq:feasible}的单点往返候选。质量和能源相关却不等价：前者限制起飞载荷，后者还受 DEM 和返程影响；体积更不能由质量替代。

Q1 不要求实体无人机分配或交付时限，故先最少架次，在同架次数内再减少能耗，最后减少累计作业时间：
\begin{equation}
\operatorname{lexmin}\left(N_1,\sum_pE_p,\sum_pt_p\right).\label{eq:q1objective}
\end{equation}
这一优先关系既响应单点运输能力问题，也避免以任意权重把“少一次飞行”与“少一点电量”相抵消。累计时间包含准备、装载、往返和交接，是作业量指标，尚未表示多机并行完工。

\subsection{锚定子集动态规划与精确性}
本题每个服务区的货箱规模允许枚举子集，且不同服务区在 Q1 中不共享调度约束，因而适合用集合划分动态规划，不需要引入更复杂的全局调度搜索。令 $D_i(U)$ 表示已交付集合 $U\subseteq\mathcal B_i$ 的最佳目标三元组，取 $D_i(\varnothing)=(0,0,0)$，其余状态初始不可达。对可行且与 $U$ 不交的候选 $S$，
\begin{equation}
D_i(U\cup S)\leftarrow\min_{\rm lex}
\{D_i(U\cup S),\ D_i(U)+(1,E_g(S),t_g(S))\}.\label{eq:dp}
\end{equation}
转移表示新增一架次恰好交付 $S$，终态 $D_i(\mathcal B_i)$ 即覆盖全部货箱且无重复的最优划分。相同 $S$ 只保留能耗最小、同能耗时时间最短的机型：它们覆盖的箱子相同，又无实体资源耦合，故删去被支配候选不影响 Q1 目标。

递推只选含最小编号未交付货箱的下一批次。每个完整划分必有一批包含该货箱，按该批先处理只是消除批次排列重复，不遗漏划分。各区分别精确求解后将目标相加，得到本问题定义下的词典序精确解。若一个服务区有 $n_i$ 箱，状态不超过 $2^{n_i}$，最坏转移规模为 $O(3^{n_i})$；精确性不意味着它能直接解决大规模实体调度。

\subsection{组批结果、约束余量与小结}
\textbf{Q1 得到 \QOneTrips{} 架次、\QOneEnergy{} kWh 和累计作业 \QOneTime{} s。}全部 \NumBoxes{} 箱恰好交付一次，最大安全载荷与完整批次见附录。图\ref{fig:batching}将箱级分配与约束占用放在一起，说明每一批为什么可行，而非只展示架次数。
\paperfig{fig04_batching}{不可拆货箱通过精确划分形成 \QOneTrips{} 个可行批次。条宽为质量、条内数字为箱数、纹理区分货类；右侧核对质量、体积和能量预算。这里展示批次，不是尚未定义的实体无人机时间表。}{batching}
验证逐批重新累加质量与体积、计算载货去程和空载返程。固定组批最低返航 SOC 为 \QOneReserveLimit\%，高于基准储备；储备改变后的重新组批见第\ref{sec:sensitivity}节。精确最优的成立范围是单点往返、给定物理规则、当前数值容差及式\eqref{eq:q1objective}，不含后续新增的实体资源、时限和通信条件。

\textbf{本问小结。}我们得到服务区—机型安全载荷边界与不可拆货箱的精确组批基准；单调性、完整子集划分和逐箱回放使其可信。但它尚不能回答哪架飞机、哪块电池何时执行。Q2 继承统一能量判据和能力认识，重新选择多点分组与排程，既不强制沿用 Q1 批次，也不能直接用单点往返载荷界代替多点路线的逐段能源检查。
\FloatBarrier
